% =====================================================================
% Appendix G --- From the input--output table to Omega: the data pipeline
% ---------------------------------------------------------------------
% Job: document, step by step, how the published Destatis input--output
%      table becomes the calibrated objects the kernel consumes --- above
%      all the input-share matrix Omega, plus alpha, lambda and the
%      consumption shares omega. Section 4 states the definitions
%      (eq:io-row, eq:io-column); Appendix C carries the A-bill leak
%      accounting; Appendix E the investment vector. This appendix fills the
%      gap between "the table" and "the shares".
% Code: OpenCGE@src/core/accounting.jl (read_data, final_demand_split,
%      generate_data, conditional_input_shares, assemble_data) and
%      OpenCGE@src/core/calibration.jl (retained_dataset for variants).
% Status: STUB --- placeholder outline (2026-10-12).
% =====================================================================
% STATUS (2026-10-23): the stub's TODO bullets are covered by the draft body below (first full draft); numbers were
%   recomputed from the raw table where possible; the import-content subsection with its two tables at the end is new.
%   Contradictions found while drafting are listed in the response letter's R2.4 cell.
\section{From the input--output table to $\Omega$: the data pipeline}\label{app:pipeline}

One published table goes in; one set of shares comes out. Every closure in the
paper is calibrated on the same objects --- the input shares $\Omega$, the
composite factor share $\alpha$, the Domar weights $\lambda$ and the household
consumption shares $\beta$ --- so differences in results cannot originate in the
data step. The accounting built on them is in Section~\ref{app:abill}.

\subsection{Source data}\label{app:pipeline-source}

The table is Germany's input--output account for 2019 (revision 2019), 71
product groups at basic prices. We read the formatted copy shipped with the
replication package, the file \texttt{I-O\_DE2019\_}
\texttt{formatiert.csv} in the \texttt{data} directory. Rows $1$--$71$ are the
modelled product groups; row $72$ (households' own goods and services) lies
outside the modelled block; rows $73$--$75$ are the domestic intermediate bill
$A$, the imported intermediates $m_i$ and the product taxes less subsidies
$t_i$, both by using sector; rows $77$--$84$ carry value added and output.
Columns $2$--$72$ are the 71 using sectors, columns $75$--$81$ the seven
final-demand categories (household consumption, government, investment,
inventories, exports). Rows and columns are located by
label, so the same code runs on the variants.

The formatting step is mechanical: a semicolon-delimited file with comma
decimals, the Destatis codes \texttt{-} and \texttt{x} read as zero, leading
blanks stripped from the label column, nothing else recoded. (A slicing trap
that shifts imports by one column, worth up to $62{,}539$~EUR~m --- nine per cent
of the $670{,}368$~EUR~m of imported intermediates --- is guarded by reading the
import row directly.)

\subsection{The intermediate-use matrix}\label{app:pipeline-z}

The $71\times71$ sector block is the intermediate-use matrix $Z$, with $x_{ui}$
the value supplied by $u$ to $i$, domestic and imported content combined; its
row and column sums reproduce eq.~\eqref{eq:io-row}--\eqref{eq:io-column} up to
the recorded imports and taxes. Those are reported by \emph{using} sector only:
row $74$ gives $m_i$ and is silent on the supplying sector. The domestic matrix $Z^{\mathrm{D}}$ therefore allocates imports
proportionally at the user level,
$Z^{\mathrm{D}}_{ui}=Z_{ui}(\text{use}_i-m_i)/\text{use}_i$: each user keeps its
recorded supplier composition, scaled by its domestic share of the bill --- a
stated assumption, not a datum.

For three users imports exceed intermediate use and the domestic remainder is
negative: coke and refined petroleum products, non-ferrous metals and water
transport --- the sectors with the highest import intensity (imports are 68, 46
and 54 per cent of their gross output). The pipeline clamps the remainder at
zero and treats them as fully import-dependent: their columns of
$Z^{\mathrm{D}}$, and hence their rows of the domestic share matrix, are zero;
nothing else is adjusted (table~\ref{tab:pipeline-clamp}).

\begin{table}[t]
\centering
\caption{Sectors whose recorded intermediate imports exceed recorded intermediate use (EUR~m). The domestic remainder is clamped at zero; the sector is treated as fully import-dependent.}
\label{tab:pipeline-clamp}
\small
\begin{tabular}{lrrrr}
\toprule
Using sector & Intermediate use & Imports & Excess & Imports / gross output \\
\midrule
Coke and refined petroleum products & $15{,}077$ & $38{,}788$ & $23{,}711$ & 0.683 \\
Non-ferrous metals (semi-finished) & $18{,}539$ & $19{,}818$ & $1{,}279$ & 0.462 \\
Water transport services & $6{,}093$ & $13{,}259$ & $7{,}166$ & 0.539 \\
\bottomrule
\end{tabular}
\end{table}

\subsection{The A-bill and the external account}\label{app:abill}

The intermediate use of a sector splits into three recorded components: the
domestic bill $A_u$, imported intermediates $M_{\mathrm{int},u}$ and product
taxes on intermediate use $T_{\mathrm{int},u}$. Together they exhaust gross
output once the value-added bill is deducted,
\begin{equation}
\sum_u \bigl( A_u + M_{\mathrm{int},u} + T_{\mathrm{int},u} \bigr)
  = \sum_u \lambda_u - 1,
\label{eq:abill}
\end{equation}
since gross output is value added plus intermediate consumption and the
$\lambda_u$ are measured in baseline income units with GDP $=1$. The kernel
charges the domestic component as input demand and books the other two as the
explicit leaks $m_i$ and $t_i$ of Section~\ref{sec:kernel}, which is what makes
the row identity \eqref{eq:io-row} close sector by sector. On the external side
the accounts satisfy
\begin{equation}
S + T_{\mathrm{int}} + M - (I + X) = F + B_{\mathrm{gov}},
\end{equation}
the identity of \eqref{eq:financing-external-account}: saving, product taxes and
imports finance investment, exports and the external position.

\subsection{Normalizing to cost shares: $\Omega$}\label{app:pipeline-omega}

Dividing each column of $Z$ by the user's intermediate bill and transposing
gives the cost shares $\Omega_{iu}=p_ux_{ui}/\sum_{u'}p_{u'}x_{u'i}$. The same operation
gives the total shares $\Omega^{\mathrm{raw}}$ from $Z$ and the domestic shares
$\Omega^{\mathrm{D}}$ from $Z^{\mathrm{D}}$, every row summing to one by
construction --- the contract the CES price index relies on. The model's
equilibrium technology is the \emph{total} matrix, which prices the intermediates
a sector actually buys; the money charged to domestic intermediates is the bill $A_i$,
with the imported and taxed content carried as the explicit leaks
$M_{\mathrm{int}}$ and $T_{\mathrm{int}}$ (Section~\ref{app:abill}). The
domestic matrix is kept for incidence and audit.

The total matrix is sparse: $1{,}631$ of its
$5{,}041$ entries (32.4 per cent) are exactly zero and half lie below $10^{-3}$.
The self-supply diagonal averages $0.224$ (median $0.201$) and is strongly
skewed --- IT and information services $0.828$, pig iron and steel $0.724$,
non-ferrous metals $0.671$, legal and management consultancy $0.569$ --- with
seven sectors above $0.45$ (13, 18, 19, 48, 53, 58, 68), the set behind the
reduced variant of Section~\ref{app:pipeline-variants}. Off the diagonal the largest
entries are water transport's purchase of warehousing and transport services
($0.543$) and insurance's from finance-related services ($0.448$). Normalization
runs down the using sector, so
columns are not normalized: summed over users, a supplier's column ranges from
$0.031$ to $4.03$ and averages $1.00$ (table~\ref{tab:pipeline-omega}).

\begin{table}[t]
\centering
\caption{Descriptive statistics of the total input-share matrix $\Omega^{\mathrm{raw}}$ (71 sectors).}
\label{tab:pipeline-omega}
\small
\begin{tabular}{lr}
\toprule
Statistic & Value \\
\midrule
Row sums (min--max) & $1.000000000000000$--$1.000000000000001$ \\
Entries exactly zero & $1{,}631$ of $5{,}041$ (32.4\,\%) \\
Entries below $10^{-3}$ & 50.0\,\% \\
Self-supply diagonal: mean / median & $0.224$ / $0.201$ \\
Self-supply diagonal: range & $0.000$--$0.828$ \\
Sectors with self-share $>0.45$ & 7 (sectors 13, 18, 19, 48, 53, 58, 68) \\
Largest off-diagonal entries & $0.543$, $0.448$, $0.407$, $0.396$, $0.395$ \\
Column sums (min--max, mean) & $0.031$--$4.03$ ($1.00$) \\
\bottomrule
\end{tabular}
\end{table}

\subsection{Value added, factor shares and Domar weights}\label{app:pipeline-va}

Value added is read in four components and reproduces gross value added to
machine precision (largest absolute difference $0.0$): compensation of employees
$1{,}763{,}874$~EUR~m, consumption of fixed capital $609{,}861$, net operating
surplus $656{,}413$ and other net taxes on production $-2{,}330$ --- gross
value added $3{,}027{,}818$~EUR~m, shares of $58.3$, $20.1$, $21.7$ and $-0.1$
per cent. The composite factor share
$\alpha_i=v_i/(p_iy_i)$ is total value added per euro of output, not the
national-accounts labour share; labour is the only primary factor and all of
$v_i$ accrues to it (Section~\ref{subsec:optimization}). It runs
from $0.043$ (coke and refined petroleum) to $0.766$ (education) and averages
$0.443$. The Domar weights $\lambda_i=y_i/\text{GDP}$ are all positive, from
$0.0002$ (fishing) to $0.1512$ (real estate services); they sum to $2.11$, and
the income weights $\lambda_i\alpha_i$ sum to exactly one.

\subsection{Final demand and household shares $\beta$}\label{app:pipeline-fd}

The seven final-demand columns form four blocks: household consumption
($1{,}309{,}448$~EUR~m), government ($648{,}260$), investment --- equipment
($240{,}444$), construction ($280{,}975$), inventories ($-29{,}801$) --- and
exports ($1{,}277{,}574$). Imported and taxed final use is reported by
category, not by supplying sector, so each category's total is split
proportionally into a domestic part and its import and product-tax components;
a non-positive category takes no domestic content.

The household shares $\beta_u$ are calibrated so that the model clears at the
baseline $y=\lambda$: the residual of
$(I-\operatorname{diag}(1-\alpha)\Omega^{\mathrm{raw}})'y$ left after government,
investment and exports is the baseline household block, clamped at zero and
renormalized to one. With the domestic-bill accounting of
Section~\ref{app:abill} that residual is exactly the observed domestic
household final demand, so no sector is genuinely clamped (one sits at
$4\times10^{-19}$ of dust). $\beta$ is concentrated --- real estate services
$0.189$, retail trade $0.088$, food $0.059$, accommodation and food services
$0.050$. The saving rate is not assumed but the data-implied residual
$s=1-\sum_u c_u/(1-\sum_u g_u)=0.1199$, with government at $0.2141$ of baseline
income, exports $0.4219$, investment $0.1624$ and intermediate imports $0.2214$
(Section~\ref{app:abill}).

\subsection{Consistency checks}\label{app:pipeline-checks}

Three identities are asserted on the real table:
$\sum_u\Omega_{iu}=1$ for every $i$, to $10^{-12}$;
$\sum_i\lambda_i\alpha_i=1$; and the intermediate-bill identity
$\sum_u A_u+\sum_i M_{\mathrm{int},i}+\sum_i T_{\mathrm{int},i}=\sum_i\lambda_i-1$,
which holds to $4\times10^{-16}$ ($\sum_uA_u=0.863$,
$\sum_iM_{\mathrm{int},i}=0.221$, $\sum_iT_{\mathrm{int},i}=0.026$,
$\sum_i\lambda_i-1=1.110$). The bill identity is pinned in the real-table block
of the paper-side test suite (\texttt{tests/test\_calibration.jl}), together with
the saving rate, the government block and the leak components; the row-sum and
income-normalization contracts are asserted for every retained-sector rebuild,
and the reproduction baseline compares every cell against the committed
\texttt{baseline/} tree to $10^{-8}$.

Production and income agree exactly, both $3{,}027{,}818$~EUR~m (difference
$0.0$). Against the expenditure side built from purchaser-price final demand
($3{,}776{,}137$~EUR~m) net of imports ($656{,}669$) and product taxes
($254{,}744$) on final use, the residual is $5.387$ per cent. That residual is a
documented open item of the raw table itself, not a calibration error of ours:
the final-demand columns are recorded at purchaser prices whose import and tax
content does not close the basic-price identity. The gross-error guard fails
above ten per cent; the measured value is well inside it
(Section~\ref{app:abill}).

\subsection{Sector-exclusion variants}\label{app:pipeline-variants}

These variants are rebuilt from the raw table, not by slicing the derived shares: the dropped sectors' rows and columns come out of
the table and the whole transformation is re-run on the slice, each retained
$\Omega$ row renormalized over the retained suppliers, $\lambda$ and $\alpha$
re-derived, and the category-level import and tax entries scaled by the retained
share. Slicing the derived matrices instead would leave every
retained user's shares summing to less than one and break the income
normalization. Dropping sector 71 leaves 70 sectors and $99.2$ per cent of gross
output; dropping also the seven high-self-share sectors leaves 63 and $87.4$ per
cent.

% TODO (AG, 2026-10-23): the import-content decomposition (direct vs embodied split, the three treatments, and the
%   closure-invariance table) was drafted here and is REMOVED for now, until the number is settled -- it is kept in
%   /tmp/appendix-imports-section.tex and summarised in the response letter's R2.4 cell. Re-insert as
%   \subsection{The import content of the programme}\label{app:pipeline-imports} when the treatment is fixed.
